\def\topic{}
\ifdefined\MyHeaderIncluded
\else
\def\qwertyflag{true}

\ifdefined\MyHeaderIncluded

\else
  \newcommand{\foobar}{FooBar}
  \def\MyHeaderIncluded{true}
% if is till the end of the file

\documentclass[10pt,a4paper]{article}
\usepackage[utf8]{inputenc}
%\usepackage[german]{babel}
%\usepackage[cyrillic]{babel}
\usepackage[T1]{fontenc}
%\usepackage[T2A]{fontenc}
%\usepackage[english, russian]{babel}
\usepackage{amsmath}
\usepackage{amsfonts}
\usepackage{amssymb}
\usepackage{amsthm}
\usepackage{mathtools} %coloneqq

\usepackage{bussproofs} %

%\renewcommand{\thesubsection}{\arabic{subsection}}
%\newtheorem{thm}{Theorem}[section]
%\renewcommand{\thethm}{\arabic{subsection}.\arabic{thm}}

\usepackage{thmtools}
\usepackage{hyperref}
\hypersetup{
    colorlinks=true,
    linkcolor=blue,
    filecolor=magenta,      
    urlcolor=cyan,
    pdftitle={Overleaf Example},
    %pdfpagemode=FullScreen,
    }
%\usepackage{stix-type1}%\turnediota
\theoremstyle{definition}

\newtheorem{constructor}{Constructor}[section]
\newtheorem{iaxiom}{Axiom}
\newtheorem{axiom}{Axiom}
\newtheorem*{axiom*}{Axiom}
\newtheorem{definition}{Definition}[section]
\newtheorem*{definition*}{Definition}
\newtheorem{remark}{Remark}[section]
\newtheorem*{remark*}{Remark}
\newtheorem{recall}{Recall}[section]
\newtheorem*{recall*}{Recall}

\newtheorem{lemma}{Lemma}[section]
\newtheorem{theorem}{Theorem}[section]
\newtheorem*{theorem*}{Theorem}
\newtheorem{note}{Note}[section]
\usepackage{etoolbox}

%\newcommand*{\@Dashv}[2]{%
%  \reflectbox{$\m@th#1#2$}%
%}
\newcommand*{\separation}[3]{% 1)binder, 2)set, 3)condition
  \{#1\in #2: #3\}
}

% https://www.freecodecamp.org/news/simple-regex-tricks-for-beginners-3acb3fa257cb/
%  \\mathcal\{T\} |---> \topology
%  \\mathcal\{([^\}]*)\}   |--->  $1
\newcommand{\AND}{\ \ \&\ \ }
\newcommand{\any}{\mathrm{\ any}}
\newcommand{\borel}{\mathcal{B}}
\newcommand{\can}{\mathrm{can}}
\newcommand{\canonize}{\mathrm{canonize}}
\newcommand{\closure}{\mbox{cl}}
\newcommand{\complex}{\mathbb{C}}
\newcommand{\converges}{\longrightarrow}
\newcommand{\countable}{\mbox{\ countable}}
\newcommand{\ldefined}{\mbox{\ defined}}
\newcommand{\dom}{\mbox{dom}}
\newcommand{\Dom}{\mbox{Dom}}
\newcommand{\Done}{\noindent \textbf{Done.}\\}
\newcommand{\lexists}{\mbox{\ exists}}
\newcommand{\lforall}{\mbox{\ forall}}
\newcommand{\lunique}{\mbox{\ unique}}
\newcommand{\family}{\mathrm{\ family}}
\newcommand{\Fm}{\mathrm{Fm}}
\newcommand{\FuncSymb}{\mathrm{FuncSymb}}
\newcommand{\function}{\mathrm{\ function}}
\newcommand{\functional}{\mbox{\ functional}}
\newcommand{\IFF}{\ \Longleftrightarrow\ }
\renewcommand{\iff}{\ \longleftrightarrow\ }
\newcommand{\IMPLIES}{\ \Longrightarrow\ }
\renewcommand{\implies}{\ \longrightarrow\ }
\newcommand{\Inf}{\mathrm{Inf}}
\newcommand{\injective}{\mbox{\ injective}}
\newcommand{\injection}{\mbox{\ injection}}
\newcommand{\interior}{\mbox{int}}
\newcommand{\intersection}{\bigcap}
\newcommand{\isempty}{\mathrm{\ empty}}
\newcommand{\LB}{\mathrm{LB}}
\newcommand{\nghbset}{\mathcal{N}}
\newcommand{\nonempty}{\mathrm{\ nonempty}}
\newcommand{\Ord}{\mathrm{Ord}}
\newcommand{\power}{\mathcal{P}}
\newcommand{\pr}{\mathrm{pr}}
\newcommand{\fst}{\mathrm{fst}}
\newcommand{\snd}{\mathrm{snd}}
\newcommand{\PredSymb}{\mathrm{PredSymb}}
\newcommand{\Proof}{\noindent \textbf{Proof.}\\}
\newcommand{\ran}{\mathrm{ran}}
\newcommand{\Ran}{\mathrm{Ran}}
\newcommand{\real}{\mathbb{R}}
\newcommand{\relation}{\mathrm{\ relation}}
\newcommand{\set}{\mathrm{\ set}}
\newcommand{\singleton}{\mathrm{\ singleton}}
\newcommand{\semantics}{\mathrm{semantics}}
\newcommand{\semialgebra}{\mathcal{S}}
\newcommand{\Sup}{\mathrm{Sup}}
\newcommand{\surjective}{\mbox{\ surjective\ on\ }}
\newcommand{\Term}{\mathrm{Term}}
\newcommand{\Type}{\mathrm{Type}}
\newcommand{\topology}{\mathcal{T}}
\newcommand{\total}{\mbox{\ total\ on\ }}
\newcommand{\transitive}{\mathrm{\ transitive}}
\newcommand{\union}{\bigcup}
\newcommand{\lTrue}{\mathrm{True}}
\newcommand{\lFalse}{\mathrm{False}}
\newcommand{\lUncertain}{\mathrm{Uncertain}}
\newcommand{\lif}{\mathrm{if}}
\newcommand{\lthen}{\mathrm{then}}
\newcommand{\lelse}{\mathrm{else}}
\newcommand{\lexception}{\mathrm{exception}}
\newcommand{\lundef}{\mathrm{undef}}
\newcommand{\univocal}{\mathrm{\ univocal}}
\newcommand{\univocalat}{\mathrm{\ univocal\ at\ }}
\newcommand{\UB}{\mathrm{UB}}
\newcommand{\Un}{\mathcal{U}}
\newcommand{\Var}{\mathrm{Var}}

\newcommand{\notvarphi}{\phi}
\renewcommand{\phi}{\varphi}

\newcommand{\comment}[1]{\quad\mbox{(#1)}}
% local commands
%\preto\section{\undef\blipo} %etoolbox
\preto\subsection{\undef\blipo} %etoolbox
% https://tex.stackexchange.com/questions/105684/limit-macro-newcommand-scope-to-section
\usepackage{graphicx}
\usepackage{authblk}

\setlength\parindent{0pt} % whole doc of \noindent
\usepackage{soul}
%\usepackage{mathgl}

\fi % guard
\ifdefined\topic
\else
\def\topic{No Titel}
\fi
\title{\topic}
\author[1]{Georgy Dunaev}
\begin{document}
\maketitle
\fi
%(Previous implementation)
\section{Relational image and glare}
In this section, relations have set arguments: $F\subseteq\V\times\V$.
The image uses an existential condition; the author's term \emph{glare}
uses the corresponding universal condition. We write $F^{-1}$ and
$F^\top$ interchangeably for the converse relation.
The remarks about families examine a different encoding and should be
read as comparisons, rather than as definitions of relational membership.
Throughout the relational calculations, the domain and range are graph
projections:
\[
\dom(F)=\cc{x}{\exists y\in\V,\ (x,y)\in F},\qquad
\ran(F)=\cc{y}{\exists x\in\V,\ (x,y)\in F}.
\]
These are not defined by testing the sethood of a totalized evaluation.
In particular, the domain of $\Sigma F$ below is its relational domain,
even when that relation is not a functional or canonical family.

\begin{definition}[(relational) membership relation]
\[
\mbox{E} \ceq \cc{(x,y)}{(y\in\V)\land(x\in y)}
\]
\end{definition}
The guard on $y$ is needed: $x\in y$ makes $x$ a set, but $y$
may be proper. Thus $\mbox{E}$ is membership restricted to two set
coordinates, agreeing with the ordinal chapter's relation $\E$.
It does not replace primitive membership of a set in an arbitrary
class. Ordinary set-pair coding is applied only to the guarded set
coordinates, without relying on any off-domain pairing convention.
\begin{remark}[Membership and family evaluation]
Membership is generally one-to-many, so it is not itself the graph of a
single-valued function. Encoding all possible outputs by a family and
then taking their union loses the individual membership relation: for a
fixed set $x$, every set $q$ belongs to some set $Y$ containing $x$ (take
$Y=\{x,q\}$). Thus the union of those possible outputs is $\V$.
\end{remark}
\begin{definition}[(relational) inverse membership relation]
\[
\invE \ceq \cc{(y,x)}{(y\in\V)\land(x\in y)} = \mbox{E}^{-1}
\]
\end{definition}
\begin{remark}[Inverse membership and family evaluation]
Likewise, the inverse membership relation lists the individual members
of a set $y$. Combining those outputs by union produces $\union y$.
This union-valued operation is a useful family evaluation, but it does
not retain the graph of the original one-to-many relation.
\end{remark}
\begin{definition}[relational image]
\[
F[A]\ceq \cc{b}{\exists a\in A,\  (a,b)\in F}
\]
\end{definition}
\begin{note}[Graph image and evaluation image]
The newer Obelisk entry \texttt{Def\allowbreak{}Relational\allowbreak{}Image} (3.26) defines an
\emph{evaluation image},
\[
\cc{b}{\exists a\in A,\ f(a)=b},
\]
using its own evaluation operation. The definition immediately above is the
\emph{graph image}, which retains each output separately. They agree
for an ordinary single-valued graph on inputs in its domain, but can
differ for a relation with several outputs.

For example, put $0=\emptyclass$, $1=\{0\}$, and
$F=\{(0,0),(0,1)\}$. The graph image of $\{0\}$ is $\{0,1\}$.
Union evaluation combines the two outputs into
$F(0)=\union\{0,1\}=1$, so its evaluation image is $\{1\}$.
The image--union and glare--intersection theorems here use the graph
convention throughout; the catalog's evaluation image is a separately
scoped construction.
\end{note}

\begin{definition}[relational glare]
\[
F\langle A \rangle\ceq \cc{b}{\forall a\in A,\  (a,b)\in F}
\]
\end{definition}
\begin{theorem}[Union as a relational image]
\[
\invE[A] = \union A
\]
\[
(\mbox{i.e. } \mbox{E}^{\top}[A]=\union A )
\]
%(or even "$\invE[A] = \union A$") % \exists
\end{theorem}
\Proof
$(\subseteq)$\\
Assume $w\in \mbox{E}^{\top}[A]$.\\
Therefore $\exists a\in A,\  (a,w)\in \mbox{E}^{\top}$.\\
So $(w,a)\in \mbox{E}$.\\
So $w\in a$.\\
Since $w\in a$ and $a\in A$, $w\in \union A$.\\
$(\supseteq)$\\
$w\in \union A$.\\
Therefore $\exists a\in A,\  w\in a$.\\
So $(w,a)\in \mbox{E}$.
So $(a,w)\in \mbox{E}^{\top}$.
From $(a,w)\in \mbox{E}^{\top}$ and $(a\in A)$, we have $(w\in \mbox{E}^{\top}[A])$.\\
\Done
\begin{theorem}[Intersection as a relational glare]
\[
\invE \langle A \rangle = \intersection A
\]
\[
(\mbox{i.e. } \mbox{E}^{\top}\langle A\rangle = \intersection A )
\]
%(or even "$\reflectbox{$\mbox{E}$}\langle A\rangle = \intersection A$")
\end{theorem}
\Proof
Fix a set $v$. This restriction supplies the sethood condition in each
comprehension, including when $A$ is empty. Then
\[
v\in \mbox{E}^{\top}\langle A \rangle
\IFF
\forall a\in A,\ (a,v)\in \mbox{E}^{\top}
\]
\[
\IFF
\forall a\in A,\ (v,a)\in \mbox{E}
\]
\[
\IFF
\forall a\in A,\ v \in a
\]
\[
\IFF
v \in \intersection A
\]
\Done
\subsection{The order structure behind glare}
The universal quantifier in glare reverses inclusion in the input.
This is a different operation from relational image, which preserves it.
The distinction yields a useful closure construction without a choice axiom.

\begin{proposition}[Image and glare under unions]
For classes $A,A',B$ and a relation $F\subseteq\V\times\V$,
\[
A\subseteq A'\implies F[A]\subseteq F[A'],
\qquad
A\subseteq A'\implies F\langle A'\rangle\subseteq F\langle A\rangle.
\]
Moreover,
\[
F[A\cup B]=F[A]\cup F[B],
\qquad
F\langle A\cup B\rangle=F\langle A\rangle\cap F\langle B\rangle.
\]
The empty-input identities are $F[\emptyclass]=\emptyclass$ and
$F\langle\emptyclass\rangle=\V$.
\end{proposition}
\Proof
An existential witness in $A$ remains a witness in $A'$. A universal
condition on $A'$ implies its restriction to $A$. Existential
quantification over a binary union is a disjunction of the two
existential conditions, while universal quantification over that union
is their conjunction. The empty cases follow from absence of witnesses
and vacuity, respectively. All resulting memberships concern sets.
\Done

\begin{proposition}[The polarity induced by a relation]
Write $p(A)=F\langle A\rangle$ and
$q(B)=F^\top\langle B\rangle$. Then
\[
B\subseteq p(A)\iff A\subseteq q(B).
\]
The operator $c(A)=q(p(A))$ is extensive, monotone, and idempotent:
\[
A\subseteq c(A),\qquad
A\subseteq A'\implies c(A)\subseteq c(A'),\qquad
c(c(A))=c(A).
\]
\end{proposition}
\Proof
Both sides of the equivalence say that $(a,b)\in F$ for every $a\in A$
and $b\in B$. Applying it to $B=p(A)$ gives extensivity. Since both
$p$ and $q$ reverse inclusion, their composite preserves it.
The symmetric argument gives $B\subseteq p(q(B))$ for every $B$.
Taking $B=p(A)$ yields $p(A)\subseteq p(q(p(A)))$. Applying $p$ to
$A\subseteq q(p(A))$ gives the reverse inclusion. Hence
$p(q(p(A)))=p(A)$, and applying $q$ proves idempotence.
\Done

\begin{example}[An order relation]
Let $P$ be a partially ordered set and
$F=\{(x,y)\in P\times P:x\leq y\}$. To keep bounds inside $P$, use
$p_P(A)=P\cap F\langle A\rangle$ and
$q_P(B)=P\cap F^\top\langle B\rangle$ for $A,B\subseteq P$.
The same proof gives a closure operator $q_Pp_P$ on subsets of $P$:
it sends a subset to the lower bounds of its upper bounds.
If $A$ has a least upper bound $s$ in $P$, this closure is exactly
$\{x\in P:x\leq s\}$. Indeed, $s$ is one of the upper bounds, so a
lower bound of all upper bounds lies below $s$; conversely, $s$ lies
below every upper bound, so every $x\leq s$ lies below them all.
The restrictions to $P$ matter for empty inputs: unrestricted glare of
an empty class is $\V$, whereas the upper bounds of the empty subset
in $P$ are precisely $P$.
\end{example}

\begin{proposition}[The right adjoint of relational image]
For a class $B$ define
\[
F^{\forall}(B)=\cc{x}{\forall y\in\V,\ (x,y)\in F\implies y\in B}.
\]
Then, for all classes $A,B$,
\[
F[A]\subseteq B\iff A\subseteq F^{\forall}(B).
\]
\end{proposition}
\Proof
The left condition says that every output paired with any $a\in A$
belongs to $B$. This is exactly the statement that each such $a$ satisfies
the defining universal condition on the right. Both implications use the
same pair witnesses, so no functionality or choice assumption is needed.
\Done

Unlike glare, this operation is monotone in $B$: enlarging the permitted
output class makes its universal condition easier to satisfy. It also
contains every input outside $\dom(F)$ by vacuity. Thus glare, direct
image, and this universal inverse operation are three distinct
constructions. For example, for the identity relation on a set $P$ and
$B\subseteq P$, one has
$F^{\forall}(B)=(\V\setminus P)\cup B$, whereas its restriction to $P$
is just $B$.

\subsection{Further identities for the relational encoding}
Here fix the expansion convention explicitly. For a relation
$F\subseteq\V\times\V$, define
\[
\Sigma F=\cc{(x,z)}{\exists y\in\V,\ (x,y)\in F\land z\in y}.
\]
Thus $(\Sigma F)[\{x\}]$ is the union of all graph outputs at $x$.
This definition uses only set arguments and outputs; it does not presume
that an absent value has any particular total-evaluation default.

\begin{theorem}[Domain decomposition]
If $F$ is functional, then
\[
\dom(\Sigma F)\cup F^\top[\{\emptyclass\}]=\dom(F),
\qquad
\dom(\Sigma F)\cap F^\top[\{\emptyclass\}]=\emptyclass.
\]
\end{theorem}
\Proof
An input belongs to $\dom(\Sigma F)$ precisely when it has a nonempty
output in $F$: an element of that output supplies the witness $z$ in
the definition. It belongs to $F^\top[\{\emptyclass\}]$ precisely when
$(x,\emptyclass)\in F$. Every output is either empty or nonempty, so the
union is $\dom(F)$, even without functionality. If an input belonged to
both pieces, functionality would identify its nonempty output with the
empty set, a contradiction. This proves disjointness.
\Done

\begin{theorem}[Recovery from nonempty values]
If $F\subseteq\V\times\V$ is functional and
$\emptyclass\notin\mbox{ran}(F)$, then
\[
F=\cc{(x,(\Sigma F)[\{x\}])}{x\in\dom(\Sigma F)}.
\]
\end{theorem}
\Proof
For $(x,y)\in F$, functionality gives
\[
z\in(\Sigma F)[\{x\}]\iff z\in y
\]
for every set $z$: any other output used in the expansion must equal
$y$. Thus the fibre is the set $y$. The nonempty-value hypothesis and
domain decomposition give $\dom(\Sigma F)=\dom(F)$. Replacing each
output by its equal fibre therefore gives exactly the displayed graph.
\Done

\begin{remark}[What expansion forgets]
For an arbitrary functional $F$, the same recovery formula is valid if
its indexing domain is replaced by $\dom(F)$. The domain is then extra
data: expansion alone forgets empty values. In particular, the empty
relation and $\{(\emptyclass,\emptyclass)\}$ both have empty expansion,
although their domains differ. Functionality also matters: the relation
$\{(0,0),(0,1)\}$, with $1=\{0\}$, has the same expansion as
$\{(0,1)\}$ and cannot be recovered from that expansion.
\end{remark}

\begin{theorem}[Image of a composite relation]
\[
(F \circ G)[X]=F[G[X]]
\]
\end{theorem}
\Proof
For a set $z$, membership in the left side means that there are $x\in X$
and a set $y$ with $(x,y)\in G$ and $(y,z)\in F$. This is exactly the
condition that $y\in G[X]$ for some $y$ with $(y,z)\in F$, or equivalently
$z\in F[G[X]]$. Class extensionality gives the equality.
\Done
%%%%%%%%%%%%%%%%%%%%%%%%%%%%%%%%%%%%%%%%%%%%%%%%
\section{A choice function and the epsilon operator}
This is a local proposal for a choice-based descriptor. Let $c$ be a
choice function on nonempty sets, and put $S_\varphi=\cc{y}{\varphi}$.
The formula below defines a chosen witness whenever $S_\varphi$ is a
nonempty set. A nonempty proper class lies outside the stated domain of
$c$; merely having a witness does not establish that $S_\varphi$ is a set.
For empty inputs, a separate partial-value or default-value convention
is needed.

\begin{definition}[Choice-based epsilon on nonempty sets]
For $S_\varphi\in\V\setminus\{\emptyclass\}$, define
\[
(\varepsilon y,\  \varphi) = c\left(\cc{y}{\varphi}\right)
\labeltagd{def}{epsilon operator for sets}
\]
Here the constant $\text{c}$ denotes a class--function of global choice(see \myref{def:globalchoice}):
\[
\text{c}:\V\setminus\{\emptyclass\}\to \V
\]
\[
\text{c}(X)\in X
\]
In the class framework with the ordinal prerequisites used in the
L\'evy chapter, limitation of size supplies such a function by the
construction below. The union and closure consequences have their
own arguments and prerequisites in the corresponding chapters; they
are not additional consequences of this definition alone.
\end{definition}

\begin{remark}[Why the original class extension needs more data]
The intended equivalence
\[
((\epsilon y,\varphi)\set)
\iff (\exists(y\set),\varphi)
\]
cannot follow merely by applying the preceding $c$ to the class
$\{y:\varphi\}$. A nonempty proper class is not an argument in the
stated domain of that class function. Moreover, an empty-set default
would itself be a set and would make the left side true even when
there is no witness. The next construction proves the intended
statement in a precise stronger setting, using a proper-class default.
It is a class-parameter operation, not an ordinary class graph with
proper classes as members.
\end{remark}

\begin{theorem}[An ordinal-ranked descriptor for arbitrary classes]
Work in the classical class setting with comprehension, $\V$ a proper
class, and a fixed injective class function $j:\V\to\On$. Assume the
usual ordinal fact that every nonempty class of ordinals has a least
member. For every class $C$, define $E_j(C)$ as follows:
\[
E_j(C)=
\begin{cases}
\text{the unique }c\in C\text{ with least }j(c),&C\ne\emptyclass,\\
\V,&C=\emptyclass.
\end{cases}
\]
This operation is well defined and extensional in $C$, and
\[
E_j(C)\set\iff C\ne\emptyclass.
\]
When $C\ne\emptyclass$, it also satisfies $E_j(C)\in C$.
Its restriction to nonempty sets is represented by an ordinary global
choice class function.
\end{theorem}
\Proof
Suppose $C$ is nonempty. Its ordinal image $j[C]$ is a nonempty
class, so it has a least member $\alpha$. Some $c\in C$ has
$j(c)=\alpha$, and injectivity of $j$ makes this $c$ unique. It is
a set because it is a member of $C$. If $C$ is empty, the assigned
value is $\V$, which is not a set by hypothesis. This proves both
the equivalence and the witness property. Equal classes have the
same ordinal images and hence the same selected member, or are both
empty, so the operation is extensional.

For the restriction, class comprehension forms the graph
\[
G=\{(A,c):A\in\V,\ A\ne\emptyclass,\ c\in A,
              \ \forall z\in A,\ j(c)\leq j(z)\}.
\]
Both coordinates here are sets, so ordinary set-pair coding suffices.
The preceding existence and uniqueness show that $G$ is a class
function on all nonempty sets and $G(A)\in A$.
\Done

The least-member fact used here does not assume that $j[C]$ is a set:
given one $\alpha_0\in j[C]$, separate
$j[C]\cap(\alpha_0+1)$ inside the set ordinal $\alpha_0+1$ and take
its least element. It is also least in the whole class. Under the
prerequisites used in the L\'evy chapter, the book's surjective-image
limitation-of-size axiom supplies the injection $j$ by choosing the
least ordinal preimage of each set. Thus the present hypothesis is a
specific established route, not an assertion that every arbitrary
choice function already acts on proper classes.

\begin{corollary}[The proposed witness equivalence, now with a setting]
For a formula $\varphi$ and its bound variable $y$, put
\[
(\epsilon_j y,\varphi):=E_j(\{y:\varphi\}).
\]
Under the preceding hypotheses,
\[
((\epsilon_j y,\varphi)\set)
\iff (\exists(y\set),\varphi).
\]
Whenever these conditions hold, the selected value satisfies
$\varphi$ under capture-avoiding substitution for $y$.
\end{corollary}
\Proof
Guarded class comprehension says that $\{y:\varphi\}$ is nonempty
exactly when there is a set witness for $\varphi$. Apply the theorem.
In that case its membership conclusion and comprehension give
$\varphi$ of the selected member. The substitution assertion uses the
ordinary total class-language semantics and its freshness conditions.
\Done

\begin{corollary}[Preserving an already chosen set selector]
If $c$ is the global choice function originally fixed on nonempty sets,
then the class-parameter operation
\[
E_{c,j}(C)=
\begin{cases}
c(C),&C\text{ is a nonempty set},\\
E_j(C),&\text{otherwise}
\end{cases}
\]
agrees with $c$ throughout its original domain and has the same
sethood and witness properties as $E_j$.
\end{corollary}
\Proof
For a nonempty set, the choice property gives a member of $C$, which
is a set. For a nonempty proper class or the empty class, the preceding
theorem gives exactly the required conclusion. These cases exhaust
the class arguments in the stated classical setting.
\Done

The ranked operator depends on the fixed $j$; a different ranking may
select a different witness. Its restriction $G$ need not equal a
previously chosen $c$, whereas the patched operation $E_{c,j}$ does
preserve that choice. Neither variation changes the proved witness
property. The empty-case value is a genuine proper
class, so this construction is not an interpretation of a descriptor
whose values are required always to be sets. It also makes no claim
that one class graph can collect pairs with arbitrary proper-class
inputs. These distinctions retain the proposed idea without identifying
it with the finite candidate semantics or with PFOL undefinedness.
\ifdefined\qwertyflag
\undef\qwertyflag
\end{document}
\fi
